\subsubsection{Lagged Recurrences as an Augmented State}

\citet{laraba2016setinvariance} review invariance for linear delay-difference
equations

\begin{equation}
    x_{t+1}=\sum_{j=0}^{L-1}A_jx_{t-j}.
\end{equation}

$A_j$ is the transition matrix multiplying the value at delay $j$. The authors
stack the current and delayed values into one state, so the delayed
equation becomes an ordinary first-order discrete-time system with a companion
transition matrix. They also define a delay-invariant set: if every initial lag
value lies in a set, every value subsequently produced by the recurrence stays
in that set. This supplies the appropriate state representation for an AST
whose leaves access several lags. In the nonlinear generator, the same idea is
used with $\mathbf{s}_t=(x_t,\ldots,x_{t-L+1})^\top$ and the set condition is
written directly as $\mathcal F_{\varepsilon}(D)\subseteq D$. The construction
extends to multivariate recurrences by stacking all $N$ variables at all $L$
lags into a state in $\mathbb R^{NL}$.
